\subsection{Grid min-marginals require a dimension-dependent budget}\label{sec:cr-star}

Let $h=2^{-j}$ and $G=h\Z\cap[0,1]$. On $[0,1]^{1+m}$ with coordinates
$(x,y_1,\ldots,y_m)$, let the uniform product grid be $G^{1+m}$, write
$U_G=\min_{G^{1+m}}F$, and, for the bag $B=\{x,y_1\}$ and $v\in G^2$,
\[
m^F(v)=\min\{F(z):z\in G^{1+m},\ (z_x,z_{y_1})=v\}.
\]
A \emph{constant-correction test} with constant $c$ retains a bag cell
$C$ of the grid if and only if $\min_{\text{corners of }C}m^F-c\le U_G$;
its coordinate version retains a grid interval $[a,a+h]$ of the centre
coordinate if and only if $\min\{V_h(a),V_h(a+h)\}-c\le U_G$, where
$V_h(a)=\min\{F(z):z\in G^{1+m},\ z_x=a\}$. On a uniform grid every node has
correction $Lh^2/8$ in every coordinate, so the corrected-grid method of
the paper applies the coordinate version with $c=(1+m)Lh^2/8$. A bag-local
budget would take $c=|B|Lh^2/8=Lh^2/4$. In both families below, the
corner minimum of each relevant bag cell equals
$\min\{V_h(a),V_h(a+h)\}$, where $[a,a+h]$ is the projection of the cell
on the centre coordinate. Hence the thresholds stated for bag cells hold
verbatim for the coordinate version.

\begin{proposition}[Constant corrections must grow with dimension]\label{prop:cr-star}
\leavevmode
\begin{enumerate}
\item[(A)] Let $j\ge1$, $m\ge1$, $\epsilon>0$ with $mh^2>4(1+\epsilon)$, and
\[
F_A(x,y)=x^2-hx\sum_{i=1}^my_i+\sum_{i=1}^my_i^2
+\Bigl(\frac{mh^2}4-1-\epsilon\Bigr)x .
\]
Its interaction graph is a star with bags $\{x,y_i\}$ ($p=2$), and $L=2$
is valid. Its unique minimizer is $x=1$, $y_i=h/2$, with $\OPT=-\epsilon$,
and any growth constant satisfies $g\le\epsilon/(1+mh^2/4)$; moreover
$U_G=0$. The only bag cell containing $(1,h/2)$ is
$C=[1-h,1]\times[0,h]$, and
\[
\min_{\text{corners of }C}m^F=(1-h)\Bigl(\frac{mh^2}4-h-\epsilon\Bigr).
\]
Hence a constant-correction test discards a cell containing the unique
optimizer unless
$c\ge(1-h)\bigl(m-4/h-4\epsilon/h^2\bigr)\,Lh^2/8$. As $m\to\infty$ the
ratio of $(1+m)Lh^2/8$ to this lower bound tends to $(1-h)^{-1}$.
\item[(B)] Let $j\ge2$, let $d\ge2/h$ be an integer, $m=d^2$, and
\[
F_B(x,y)=\Bigl(x-\frac12\Bigr)^2+\sum_{i=1}^m
\Bigl(y_i-\frac h2-\frac{x-1/2}{d}\Bigr)^2 .
\]
It is a star with $p=2$; $L=4$ is valid; its unique minimizer
$z^*=(1/2,h/2,\ldots,h/2)$ is interior; and
$F_B(z)-\OPT\ge\frac{3-\sqrt5}2\norm{z-z^*}^2$, so
$\kappa\le2(3+\sqrt5)<11$ for every $d$ and $h$. Moreover
$U_G=\frac12-\frac{dh}2+\frac{d^2h^2}4$, and both bag cells containing
$(1/2,h/2)$ have
\[
\min_{\text{corners}}m^F-U_G=dh\Bigl(\frac12-h\Bigr)+2h^2-\frac12 .
\]
Thus a constant-correction test discards every cell containing the
unique optimizer unless $c\ge dh(\frac12-h)+2h^2-\frac12$, a quantity of
order $\sqrt m\,h$ at fixed $h$, while $p=2$ and $\kappa<11$.
\end{enumerate}
\end{proposition}
\begin{proof}
(A) Write $\beta=mh^2/4-1-\epsilon>0$. All coordinate second derivatives
equal $2$ and $F_A$ is quadratic, so $L=2$ is valid. For fixed $x$, each
leaf term $y^2-hxy$ is strictly convex with minimizer $hx/2\in[0,1]$ and
minimum $-h^2x^2/4$. Hence $\min_yF_A(x,y)=(1-mh^2/4)x^2+\beta x=:V(x)$,
which is strictly concave because $mh^2>4$, with $V(0)=0$ and
$V(1)=-\epsilon$. For $x\in(0,1)$, strict concavity gives
$V(x)>(1-x)V(0)+xV(1)\ge-\epsilon$. So $x=1$ is the unique minimizing
centre and $y_i=h/2$ the unique leaves; $\OPT=-\epsilon$. The origin has
$F_A=0=\OPT+\epsilon$ and squared distance $1+mh^2/4$ from the minimizer,
which bounds $g$.

On the grid, $y(y-hx)\ge0$ for $y\in G$ and $x\in[0,1]$, since either
$y=0$ or $y\ge h\ge hx$. Also $x^2+\beta x\ge0$ for $x\ge0$. Hence
$F_A\ge0$ on $G^{1+m}$ with equality at the origin, so $U_G=0$, and every
leaf's grid minimum is $0$. The point $(1,h/2)$ has $x$ at the right
endpoint and $h/2$ interior to $[0,h]$, so $C$ is the only grid cell of
the bag containing it. At its corners,
$m^F(x,y_1)=x^2+\beta x+(y_1^2-hxy_1)$. For $x=1-h$ this equals
$(1-h)(1-h+\beta)=(1-h)(mh^2/4-h-\epsilon)$ at $y_1=0$ and exceeds it by
$h^3$ at $y_1=h$. For $x=1$ it equals $mh^2/4-\epsilon$ at both $y_1=0,h$,
which exceeds the previous value by $h(mh^2/4-\epsilon)+h(1-h)>0$. This
proves the corner minimum; it also equals $\min\{V_h(1-h),V_h(1)\}$,
because every leaf's grid minimum is attained at $0$. Since $U_G=0$ and
$L=2$, the threshold for $c$ is as stated, and
$(1+m)h^2/4\big/\bigl[(1-h)(mh^2/4-h-\epsilon)\bigr]\to(1-h)^{-1}$.

(B) The Hessian $H$ has $H_{xx}=2+2m/d^2=4$, $H_{y_iy_i}=2$,
$H_{xy_i}=-2/d$, and zero elsewhere; thus $L=4$ is valid. Vectors with
zero $x$-component and leaf components summing to zero are eigenvectors
with eigenvalue $2$. On the span of $e_x$ and $m^{-1/2}\mathbf 1_y$, $H$
acts as $\bigl(\begin{smallmatrix}4&-2\\-2&2\end{smallmatrix}\bigr)$, with
eigenvalues $3\pm\sqrt5$. Since $F_B\ge0$ vanishes only at $z^*\in(0,1)^{1+m}$
and $F_B(z)=\frac12(z-z^*)^{\T}H(z-z^*)$, growth holds with
$g=(3-\sqrt5)/2$ and $\kappa=L/g=2(3+\sqrt5)$.

Write $t=x-1/2\in h\Z\cap[-1/2,1/2]$ for centre grid values (here
$1/2\in G$ because $j\ge1$). A leaf's ideal value
$\iota(t)=h/2+t/d$ lies in $[0,h]$ because $|t|/d\le1/(2d)\le h/4$.
Its distance to $G$ is therefore $h/2-|t|/d$, and the conditional grid
value of the centre is
\[
V_h(t)=t^2+m\Bigl(\frac h2-\frac{|t|}d\Bigr)^2=2t^2-dh|t|+\frac{d^2h^2}4 .
\]
On the grid $|t|\le1/2\le dh/4$, and $s\mapsto2s^2-dhs$ decreases on
$[0,dh/4]$, so $U_G=V_h(\pm1/2)=\frac12-\frac{dh}2+\frac{d^2h^2}4$.
The optimizer's bag projection $(1/2,h/2)$ lies on the common edge of
the cells $[1/2-h,1/2]\times[0,h]$ and $[1/2,1/2+h]\times[0,h]$ and in no
other cell. At $t=0$ both leaf labels $0,h$ are at distance $h/2$ from
$\iota$, so $m^F=d^2h^2/4$. At $t=\pm h$ the nearer label gives
$m^F=V_h(h)=2h^2-dh^2+d^2h^2/4$ and the farther adds $2h^2/d$. Since
$d\ge2$, each cell has corner minimum $d^2h^2/4-(d-2)h^2
=\min\{V_h(0),V_h(\pm h)\}$, and subtracting $U_G$ gives the stated
difference. For $h\le1/4$ this difference is at least
$dh/4+2h^2-1/2$, which grows linearly in $d=\sqrt m$. Discarding both
centre intervals adjacent to $x=1/2$ removes the optimizer.
\end{proof}

\begin{example}[The smallest instance]\label{ex:cr-star32}
Take $h=1/2$, $m=32$, $\epsilon=1/16$ in (A):
$F=x^2-\frac x2\sum y_i+\sum y_i^2+\frac{15}{16}x$ on $[0,1]^{33}$.
The exact leaf recourse is $y_i=x/4$, with $V(x)=-x^2+\frac{15}{16}x$, so
the unique optimizer is $x=1$, $y_i=1/4$, value $-1/16$. On
$G=\{0,\frac12,1\}$ the conditional grid values at $x=0,\frac12,1$ are
$0,\frac{23}{32},\frac{31}{16}$, and $U_G=0$. The four corner
min-marginals of $C=[\frac12,1]\times[0,\frac12]$ are
$\frac{23}{32},\frac{27}{32},\frac{31}{16},\frac{31}{16}$. The bag-local
budget $\frac18$ rejects $C$; the global correction $\frac{33}{16}$ keeps it.
The conditional grid error is
\[
V_h(x)-V(x)=m\,\dist(x/4,G)^2,
\]
which varies from $0$ at $x=0$ to $m/16$ at $x=1$; no normalization of
messages removes this variation. Because the growth constant of family (A)
degrades as $m$ grows, this instance refutes budgets depending on
$(p,L)$ only; family (B) refutes budgets depending on $(p,\kappa)$.
\end{example}

